\section{Resultados}

\subsection{Actividad 1}

Tras la compilación y ejecución del programa, el resultado final fue una ventana gráfica capaz de alterar su color de fondo en tiempo real al presionar las teclas asignadas. Las pruebas en pantalla demostraron que la transición de colores ocurre de manera inmediata, logrando que el repintado del búfer coincida exactamente con los tonos de la paleta solicitada en las instrucciones de la práctica. Cabe destacar que las capturas de pantalla con los resultados finales en funcionamiento ya fueron presentadas y referenciadas en la sección de Experimentos.

Al evaluar este ejercicio, y cumpliendo con el análisis de la metodología requerida, se observaron los siguientes puntos sobre nuestro enfoque:

\begin{itemize}
    \item \textbf{Ventajas del método:} La principal ventaja de realizar la conversión matemática a lápiz y alimentar a \texttt{glClearColor} directamente con valores flotantes normalizados es la optimización del rendimiento. El hardware de la tarjeta gráfica está diseñado para procesar números de punto flotante de manera nativa. Al entregarle los datos ya digeridos entre 0.0f y 1.0f, evitamos que el procesador desperdicie ciclos de reloj traduciendo cadenas de texto hexadecimales o enteros de 8 bits en cada iteración del bucle principal de renderizado.
    \item \textbf{Desventajas del método:} El punto débil de esta metodología es su rigidez a nivel de ingeniería de software. Al dejar los valores flotantes "quemados" (hardcoded) dentro de los condicionales de la función \texttt{keyboardCallback}, el código pierde escalabilidad. Si en el futuro se requiriera implementar una interfaz para que el usuario escoja un color personalizado, nuestro enfoque actual nos obligaría a recalcular a mano, alterar el código fuente y recompilar todo el proyecto. Una solución más elegante habría sido desarrollar una función auxiliar en C++ que parsee automáticamente un string hexadecimal (ej. "\#00C2CB") y retorne un vector normalizado en tiempo de ejecución.
\end{itemize}

\subsection{Actividad 2}

[PENDIENTE: integrar resultados de la actividad 2 proporcionados por el integrante responsable]

\subsection{Actividad 3}

Se obtuvo un semicírculo relleno centrado en el origen, cuyo resultado en pantalla se muestra en la Figura \ref{fig:mediocirculo} de la sección de Experimentos. El trazo aparece cerrado sobre el diámetro y sin discontinuidades en el arco, lo que confirma que el recorrido angular limitado a $180^\circ$ genera la mitad exacta de la circunferencia.\\

El resultado más relevante de esta actividad fue comprobar que una misma lista de vértices produce figuras distintas según la primitiva con que se interprete. Con \texttt{GL\_TRIANGLES} la geometría se desplegaba fragmentada, ya que los vértices se agrupan de tres en tres sin considerar el vértice central; al emplear \texttt{GL\_TRIANGLE\_FAN} ese primer vértice pasa a funcionar como pivote común y la superficie se rellena por completo. La corrección no requirió modificar una sola coordenada, únicamente el criterio de conexión entre los vértices ya existentes.\\

Cabe señalar que la resolución de la figura queda determinada por el incremento angular. Al conservarse el paso de $15^\circ$, el contorno se aproxima mediante doce segmentos rectos, lo cual resulta perceptible en el borde curvo; reducir dicho incremento suavizaría el arco a costa de generar un mayor número de vértices.\\

\subsection{Actividad 4}

Se obtuvo un pentágono regular relleno, centrado en el origen y con los cinco colores de la paleta solicitada asignados uno por vértice. La conversión de los valores hexadecimales al formato normalizado se verificó de manera visual comparando el tono de cada vértice contra la paleta del manual. El degradado observado en el interior de la figura confirma que la interpolación de los atributos de color ocurre en la etapa de rasterización, entre el \textit{vertex shader} y el \textit{fragment shader}, sin necesidad de declarar color alguno por fragmento.\\

La variación mostró que la representación de una misma geometría depende del arreglo de índices: al suprimirlo, los cinco vértices se despliegan como una polilínea cerrada en lugar de tres triángulos, sin modificar una sola coordenada.\\

\subsection{Actividad 5}

La gráfica de $\tan(2\theta)$ se desplegó completa sobre el dominio $[-2\pi : 2\pi]$, con sus nueve ramas continuas separadas por las ocho asíntotas verticales previstas analíticamente en $\theta = \pi/4 + k\pi/2$. El coloreado por rama permitió comprobar que la segmentación obtenida coincide con la esperada.\\

El recorte de las muestras que exceden el intervalo visible resultó indispensable, ya que sin él los vértices cercanos a las asíntotas se proyectan a valores arbitrariamente grandes. La variación del factor de escala evidenció que dicho parámetro no altera la función, sino la porción de cada rama que alcanza a mostrarse antes de salir del área de dibujo.\\
